\chapter{LightGBM: A Highly Efficient Gradient Boosting Decision Tree}

\begin{note}
完整英文原文见本地 PDF：
\texttt{papers/lightgbm\_a\_highly\_efficient\_gradient\_boosting\_decision\_tree.pdf}。
本章保留原论文标题、章节结构、核心公式、算法与图表位置，便于和 PDF 原文逐节对照。
\end{note}

\section{Abstract}

\textbf{原文位置：}PDF p.1，Abstract。

\textbf{结构要点：}
\begin{itemize}
  \item GBDT is popular and effective, with implementations such as XGBoost and pGBRT.
  \item Existing implementations are inefficient when feature dimension and data size are large.
  \item The paper proposes Gradient-based One-Side Sampling (GOSS).
  \item The paper proposes Exclusive Feature Bundling (EFB).
  \item LightGBM combines GOSS and EFB and can speed up conventional GBDT by over 20 times with almost the same accuracy.
\end{itemize}

\section{1 Introduction}

\textbf{原文位置：}PDF p.1--p.2，Section 1。

\textbf{结构要点：}
\begin{itemize}
  \item GBDT is widely used because of efficiency, accuracy and interpretability.
  \item Big data brings new challenges in both number of features and number of instances.
  \item Conventional GBDT scans all data instances for all possible split points, so complexity depends on both data size and feature size.
  \item The natural acceleration direction is to reduce data instances and reduce features.
  \item GOSS reduces data instances.
  \item EFB reduces effective feature number.
  \item The new algorithm with GOSS and EFB is called LightGBM.
\end{itemize}

\section{2 Preliminaries}

\subsection{2.1 GBDT and Its Complexity Analysis}

\textbf{原文位置：}PDF p.2，Section 2.1，Algorithm 1。

\textbf{结构要点：}
\begin{itemize}
  \item GBDT is an ensemble model of decision trees trained in sequence.
  \item Each iteration fits negative gradients, also known as residual errors.
  \item The main cost is learning decision trees.
  \item The most time-consuming part is finding best split points.
  \item Histogram-based algorithm buckets continuous feature values into discrete bins.
  \item Histogram building costs $O(\#data\times \#feature)$.
  \item Split point finding costs $O(\#bin\times \#feature)$.
  \item Since $\#bin$ is much smaller than $\#data$, histogram building dominates computational complexity.
\end{itemize}

\subsection{2.2 Related Work}

\textbf{原文位置：}PDF p.2--p.3，Section 2.2。

\textbf{结构要点：}
\begin{itemize}
  \item Existing GBDT implementations include XGBoost, pGBRT, scikit-learn and R gbm.
  \item XGBoost supports both pre-sorted and histogram-based algorithms.
  \item Random data sampling can hurt accuracy in GBDT.
  \item Feature filtering depends on redundancy assumptions and may hurt accuracy.
  \item Sparse data are common in large-scale real applications.
  \item Histogram-based GBDT lacks efficient sparse optimization in previous systems.
  \item LightGBM addresses these limits with GOSS and EFB.
\end{itemize}

\section{3 Gradient-based One-Side Sampling}

\subsection{3.1 Algorithm Description}

\textbf{原文位置：}PDF p.3--p.4，Section 3.1，Algorithm 2。

\textbf{结构要点：}
\begin{itemize}
  \item AdaBoost has sample weights, but GBDT has no native sample weights.
  \item Gradients in GBDT provide useful information for sampling.
  \item Small-gradient instances are already well trained.
  \item Simply dropping small-gradient instances changes data distribution.
  \item GOSS keeps all large-gradient instances and randomly samples small-gradient instances.
  \item Sampled small-gradient instances are multiplied by $\frac{1-a}{b}$ when computing information gain.
\end{itemize}

\subsection{3.2 Theoretical Analysis}

\textbf{原文位置：}PDF p.4--p.5，Section 3.2，Definition 3.1，Theorem 3.2。

\textbf{原文公式：}
\[
  \tilde{V}_j(d)
  =
  \frac{1}{n}
  \left[
    \frac{
      \left(
        \sum_{x_i\in A_l}g_i
        +
        \frac{1-a}{b}\sum_{x_i\in B_l}g_i
      \right)^2
    }{n_l^j(d)}
    +
    \frac{
      \left(
        \sum_{x_i\in A_r}g_i
        +
        \frac{1-a}{b}\sum_{x_i\in B_r}g_i
      \right)^2
    }{n_r^j(d)}
  \right].
\]

\textbf{结构要点：}
\begin{itemize}
  \item GOSS uses an estimated variance gain $\tilde{V}_j(d)$ on subset $A\cup B$.
  \item $A$ contains large-gradient instances.
  \item $B$ contains sampled small-gradient instances.
  \item The coefficient $\frac{1-a}{b}$ normalizes the small-gradient subset back to the scale of the original small-gradient data.
  \item Theorem 3.2 bounds the approximation error of GOSS.
  \item Random sampling is a special case of GOSS when $a=0$.
  \item GOSS can outperform random sampling under the same sampling ratio.
\end{itemize}

\section{4 Exclusive Feature Bundling}

\textbf{原文位置：}PDF p.5--p.6，Section 4，Algorithm 3，Algorithm 4，Theorem 4.1。

\textbf{结构要点：}
\begin{itemize}
  \item High-dimensional data are usually sparse.
  \item Many sparse features are mutually exclusive.
  \item Mutually exclusive features rarely take nonzero values simultaneously.
  \item Exclusive features can be bundled into a single feature bundle.
  \item Histogram building complexity changes from $O(\#data\times \#feature)$ to $O(\#data\times \#bundle)$.
  \item Finding the optimal bundling strategy is NP-hard.
  \item LightGBM reduces the bundling problem to graph coloring and uses a greedy algorithm.
  \item Features in one bundle are merged by assigning different bin ranges through offsets.
\end{itemize}

\section{5 Experiments}

\textbf{原文位置：}PDF p.6--p.8，Section 5，Table 1，Table 2，Table 3，Figure 1，Figure 2。

\textbf{结构要点：}
\begin{itemize}
  \item Experiments use five public datasets: Allstate, Flight Delay, LETOR, KDD10 and KDD12.
  \item Datasets cover sparse features, dense features, binary classification and ranking.
  \item Baselines include XGBoost and LightGBM without GOSS/EFB.
  \item LightGBM is the fastest while maintaining almost the same accuracy.
  \item Compared with lgb\_baseline, LightGBM speeds up training by 21x, 6x, 1.6x, 14x and 13x on five datasets.
  \item XGBoost histogram version runs out of memory on KDD10 and KDD12.
  \item GOSS and EFB do not hurt accuracy while bringing speed-up.
\end{itemize}

\subsection{5.1 Overall Comparison}

\textbf{原文位置：}PDF p.6--p.8，Section 5.1。

\textbf{结构要点：}
\begin{itemize}
  \item XGBoost and lgb\_baseline are used as baselines.
  \item xgb\_exa uses pre-sorted algorithm.
  \item xgb\_his uses histogram-based algorithm.
  \item LightGBM combines lgb\_baseline with GOSS and EFB.
  \item LightGBM achieves large training speed-up with nearly unchanged accuracy.
\end{itemize}

\subsection{5.2 Analysis on GOSS}

\textbf{原文位置：}PDF p.8，Section 5.2，Table 4。

\textbf{结构要点：}
\begin{itemize}
  \item GOSS alone can bring nearly 2x speed-up using 10\%--20\% data.
  \item GOSS still needs full-data prediction and gradient computation, so speed-up is not linear in sampling ratio.
  \item GOSS is compared with Stochastic Gradient Boosting (SGB).
  \item Under the same sampling ratio, GOSS is consistently more accurate than SGB.
\end{itemize}

\subsection{5.3 Analysis on EFB}

\textbf{原文位置：}PDF p.8，Section 5.3。

\textbf{结构要点：}
\begin{itemize}
  \item EFB contribution is measured by comparing lgb\_baseline and EFB\_only.
  \item EFB gives significant speed-up on large-scale sparse datasets.
  \item EFB merges one-hot and implicitly exclusive sparse features into fewer features.
  \item Bundling improves spatial locality and cache hit rate.
\end{itemize}

\section{6 Conclusion}

\textbf{原文位置：}PDF p.8，Section 6。

\textbf{结构要点：}
\begin{itemize}
  \item LightGBM is a new GBDT algorithm with GOSS and EFB.
  \item GOSS handles large numbers of data instances.
  \item EFB handles large numbers of features.
  \item Theory and experiments show that LightGBM improves speed and memory consumption while maintaining accuracy.
  \item Future work includes choosing optimal $a$ and $b$ in GOSS and improving EFB.
\end{itemize}

\chapter{《LightGBM: A Highly Efficient Gradient Boosting Decision Tree》中文版}

\section{摘要}

梯度提升决策树（GBDT）是一种流行的机器学习算法，并且已经有 XGBoost 和 pGBRT 等有效实现。
尽管这些实现采用了许多工程优化，但当特征维度很高、数据规模很大时，它们的效率和可扩展性仍然不够理想。
主要原因是：对于每个特征，算法都需要扫描所有数据样本，以估计所有可能切分点的信息增益，这非常耗时。

为了解决这个问题，论文提出了两种新技术：基于梯度的单边采样
（Gradient-based One-Side Sampling, GOSS）和互斥特征捆绑
（Exclusive Feature Bundling, EFB）。GOSS 会排除大量小梯度样本，
只使用剩余样本估计信息增益。由于大梯度样本在信息增益计算中更重要，
GOSS 可以用更小的数据量获得较准确的信息增益估计。EFB 将互斥特征捆绑在一起，
从而减少特征数量。结合 GOSS 和 EFB 的新 GBDT 实现被称为 LightGBM。
实验表明，LightGBM 在几乎保持相同精度的同时，可以使传统 GBDT 训练速度提升超过 20 倍。

\section{1 引言}

GBDT 因效率、精度和可解释性而被广泛使用，在多分类、点击预测、学习排序等任务中表现突出。
随着数据规模和特征数量不断增大，GBDT 面临精度与效率之间的新挑战。
传统 GBDT 实现通常需要对每个特征扫描所有样本，以估计所有候选切分点的信息增益。
因此，其计算复杂度同时依赖样本数量和特征数量，当数据很大时会非常耗时。

一个直接的加速思路是减少样本数量和特征数量，但这并不简单。
对于样本，GBDT 没有像 AdaBoost 那样的天然样本权重，因此不能直接使用 AdaBoost 的采样方法。
LightGBM 观察到，梯度大小可以反映样本对信息增益的重要性，于是提出 GOSS：
保留大梯度样本，只随机丢弃小梯度样本。对于特征，实际高维特征空间通常很稀疏，
很多特征几乎不会同时非零，因此可以将这些互斥特征捆绑成更少的特征，这就是 EFB。

\section{2 预备知识}

\subsection{2.1 GBDT 及其复杂度分析}

GBDT 是由按顺序训练的多棵决策树组成的集成模型。每次迭代中，GBDT 通过拟合损失函数的负梯度
来训练新的决策树。GBDT 的主要计算成本在于学习决策树，而学习决策树最耗时的部分是寻找最佳分裂点。

预排序算法会在排好序的特征值上枚举所有可能分裂点，虽然可以找到最优分裂点，但训练速度和内存消耗都较差。
直方图算法将连续特征值离散为若干桶，并在训练中构建特征直方图。该算法在内存和速度上更高效，
因此 LightGBM 在直方图算法基础上展开。

直方图构建的复杂度为 $O(\#data\times \#feature)$，
寻找分裂点的复杂度为 $O(\#bin\times \#feature)$。
由于桶数量通常远小于样本数量，直方图构建主导了整体计算复杂度。
因此，如果能减少样本数或特征数，就能显著加速 GBDT 训练。

\subsection{2.2 相关工作}

已有 GBDT 实现包括 XGBoost、pGBRT、scikit-learn 和 R 中的 gbm。
scikit-learn 和 R gbm 使用预排序算法，pGBRT 使用直方图算法，XGBoost 同时支持预排序和直方图算法。
为了减少训练数据规模，常见方法是下采样样本，但随机采样可能损害 GBDT 精度。
为了减少特征数量，也可以过滤弱特征，但这依赖特征冗余假设，实践中不一定成立。

实际大规模数据通常非常稀疏。预排序算法可以通过忽略零值降低训练成本，
但直方图算法之前缺少有效的稀疏优化方案。LightGBM 正是为了解决这些限制，提出了 GOSS 和 EFB。

\section{3 基于梯度的单边采样}

\subsection{3.1 算法描述}

在 AdaBoost 中，样本权重可以表示样本重要性。但 GBDT 没有天然样本权重，
因此 AdaBoost 的采样方法不能直接使用。LightGBM 注意到，GBDT 中每个样本的梯度提供了有用的采样信息。
如果一个样本梯度很小，说明该样本训练误差较小，已经被模型学得比较好。

直接丢弃小梯度样本会改变数据分布，从而损害模型精度。GOSS 的做法是：
保留所有大梯度样本，并从小梯度样本中随机采样一部分。为了补偿数据分布变化，
在计算信息增益时，对被采样的小梯度样本乘以常数 $\frac{1-a}{b}$。
这样既能更关注欠训练样本，又不会过度改变原始数据分布。

\subsection{3.2 理论分析}

GBDT 使用决策树从输入空间学习到梯度空间的函数。对某个固定节点，特征 $j$ 在切分点 $d$ 上的信息增益
可以用分裂后的方差增益表示。GOSS 将样本按梯度绝对值降序排列，保留前 $a\times 100\%$ 的大梯度样本形成集合 $A$，
再从剩余小梯度样本中随机采样比例为 $b$ 的集合 $B$。

GOSS 在 $A\cup B$ 上估计方差增益：
\[
  \tilde{V}_j(d)
  =
  \frac{1}{n}
  \left[
    \frac{
      \left(
        \sum_{x_i\in A_l}g_i
        +
        \frac{1-a}{b}\sum_{x_i\in B_l}g_i
      \right)^2
    }{n_l^j(d)}
    +
    \frac{
      \left(
        \sum_{x_i\in A_r}g_i
        +
        \frac{1-a}{b}\sum_{x_i\in B_r}g_i
      \right)^2
    }{n_r^j(d)}
  \right].
\]

其中 $\frac{1-a}{b}$ 用于把小梯度采样集合的梯度和归一化回原始小梯度数据规模。
论文的定理表明，GOSS 的近似误差可以被控制，并且随机采样是 $a=0$ 时的特殊情形。
在许多情况下，GOSS 在相同采样率下可以优于随机采样。

\section{4 互斥特征捆绑}

高维数据通常非常稀疏，这为近似无损地减少特征数量提供了可能。
在稀疏特征空间中，很多特征是互斥的，也就是说它们几乎不会同时取非零值。
这些互斥特征可以被安全地捆绑成一个特征，称为互斥特征束。

通过精心设计的特征扫描算法，可以从特征束中构建出与单独特征相同的直方图。
这样，直方图构建复杂度可以从
$O(\#data\times \#feature)$ 降为
$O(\#data\times \#bundle)$，其中 $\#bundle$ 远小于 $\#feature$。

EFB 需要解决两个问题：第一，哪些特征应该捆绑在一起；第二，如何构造捆绑后的特征。
论文证明，寻找最少互斥特征束的最优划分是 NP-hard 问题。
因此，LightGBM 将该问题转化为图着色问题：把特征看作顶点，如果两个特征不互斥，就在它们之间连边；
然后使用贪婪算法得到近似解。

对于特征合并，关键是要保证原始特征值仍能从特征束中识别出来。
由于直方图算法存储的是离散桶，LightGBM 可以通过给不同特征的桶添加偏移量，
让多个互斥特征占据不同的桶区间，从而安全地合并成一个特征束。

\section{5 实验}

论文使用五个公开数据集：Allstate、Flight Delay、LETOR、KDD10 和 KDD12。
这些数据集覆盖稀疏特征、稠密特征、二分类和排序任务。
实验环境是一台 Linux 服务器，使用两颗 E5-2670 v3 CPU，共 24 核和 256GB 内存。

\subsection{5.1 整体比较}

实验使用 XGBoost 和不含 GOSS/EFB 的 LightGBM 版本作为基线。
XGBoost 分为预排序版本 xgb\_exa 和直方图版本 xgb\_his。
LightGBM 则是在 lgb\_baseline 基础上加入 GOSS 和 EFB。

结果表明，LightGBM 在几乎保持相同测试精度的同时训练最快。
与 lgb\_baseline 相比，LightGBM 在 Allstate、Flight Delay、LETOR、KDD10、KDD12
上的训练加速比分别为 21 倍、6 倍、1.6 倍、14 倍和 13 倍。
XGBoost 的直方图版本在 KDD10 和 KDD12 上由于内存不足无法运行。
这些结果说明，GOSS 和 EFB 在带来显著加速的同时不会损害精度。

\subsection{5.2 GOSS 分析}

通过比较 LightGBM 和 EFB\_only，论文分析 GOSS 的加速能力。
结果显示，在只使用 10\%--20\% 数据的情况下，GOSS 本身可以带来接近 2 倍的加速。
不过 GOSS 仍然需要在全量数据上进行预测和计算梯度，因此整体加速并不会与采样比例完全线性相关。

论文还将 GOSS 与随机梯度提升（SGB）比较。在相同采样率下，GOSS 的精度始终高于 SGB。
这说明 GOSS 是比普通随机采样更有效的采样方法。

\subsection{5.3 EFB 分析}

论文通过比较 lgb\_baseline 和 EFB\_only 来分析 EFB 的贡献。
结果显示，EFB 在大规模稀疏数据集上可以带来显著加速。
这是因为 EFB 将大量 one-hot 特征和隐式互斥稀疏特征合并成更少的特征。
此外，由于许多原本分散的特征被捆绑在一起，空间局部性和缓存命中率也会提升。
因此，EFB 是直方图算法中利用稀疏性的有效方法。

\section{6 结论}

本文提出了一种新的 GBDT 算法 LightGBM，其中包含两项新技术：
基于梯度的单边采样（GOSS）和互斥特征捆绑（EFB）。
GOSS 用于处理大量样本，EFB 用于处理大量特征。
论文对两项技术进行了理论分析和实验验证。
实验结果与理论分析一致，表明在 GOSS 和 EFB 的帮助下，LightGBM 可以在计算速度和内存消耗方面显著优于
XGBoost 和 SGB，同时保持相近精度。未来工作包括研究 GOSS 中 $a$ 和 $b$ 的最优选择，
以及继续提升 EFB 在更多特征场景下的表现。

